\documentclass[12pt]{article}

\usepackage[utf8]{inputenc}
\usepackage[spanish]{babel}
\usepackage[shortlabels]{enumitem}

\usepackage{mathtools}
\usepackage{amssymb}
\usepackage{amsthm}

\usepackage{geometry}
\geometry{margin=2.5cm}

\usepackage{hyperref}

\title{Soluciones Ejercicios Aplicaciones Lineales}
\author{}

\begin{document}

\section*{Solución Ejercicio 1}

Tenemos el endomorfismo
\[
f:\mathbb{R}^3\longrightarrow \mathbb{R}^3
\]
definido por
\[
f(x_1,x_2,x_3)=(x_1+x_2+x_3,\;x_1+x_2-x_3,\;x_3).
\]

La matriz asociada a \(f\) en la base canónica es
\[
M_C(f)=
\begin{pmatrix}
1&1&1\\
1&1&-1\\
0&0&1
\end{pmatrix}.
\]

\subsubsection*{a) Núcleo e imagen de la aplicación lineal}

Calculamos primero el núcleo. Para ello resolvemos:
\[
f(x_1,x_2,x_3)=(0,0,0).
\]

Es decir:
\[
(x_1+x_2+x_3,\;x_1+x_2-x_3,\;x_3)=(0,0,0).
\]

Por tanto:
\[
\begin{cases}
x_1+x_2+x_3=0,\\
x_1+x_2-x_3=0,\\
x_3=0.
\end{cases}
\]

De la tercera ecuación:
\[
x_3=0.
\]

Sustituyendo en la primera:
\[
x_1+x_2=0.
\]

Por tanto:
\[
x_2=-x_1.
\]

Así, los vectores del núcleo son:
\[
(x_1,x_2,x_3)=(x_1,-x_1,0).
\]

Luego:
\[
(x_1,x_2,x_3)=x_1(1,-1,0).
\]

Por tanto,
\[
\ker f=L<\{(1,-1,0)\}>.
\]

Una base de \(\ker f\) es
\[
\boxed{\{(1,-1,0)\}.}
\]

Luego:
\[
\boxed{\dim(\ker f)=1.}
\]

Ahora calculamos la imagen. A partir de la matriz:
\[
M_C(f)=
\begin{pmatrix}
1&1&1\\
1&1&-1\\
0&0&1
\end{pmatrix},
\]
la imagen está generada por sus columnas:
\[
\operatorname{Im} f
=
L<\{(1,1,0),(1,1,0),(1,-1,1)\}>.
\]

Como la primera y la segunda columna son iguales, tenemos:
\[
\operatorname{Im} f
=
L<\{(1,1,0),(1,-1,1)\}>.
\]

Comprobamos que estos dos vectores son linealmente independientes. Si
\[
\alpha(1,1,0)+\beta(1,-1,1)=(0,0,0),
\]
entonces:
\[
(\alpha+\beta,\alpha-\beta,\beta)=(0,0,0).
\]

De la tercera coordenada:
\[
\beta=0.
\]

Sustituyendo:
\[
\alpha+\beta=0
\Rightarrow
\alpha=0.
\]

Por tanto, los vectores son linealmente independientes.

Así, una base de la imagen es
\[
\boxed{\{(1,1,0),(1,-1,1)\}.}
\]

Luego:
\[
\boxed{\dim(\operatorname{Im} f)=2.}
\]

También podemos expresar la imagen en forma implícita.

Sea
\[
(y_1,y_2,y_3)\in \operatorname{Im} f.
\]

Entonces:
\[
(y_1,y_2,y_3)=\alpha(1,1,0)+\beta(1,-1,1).
\]

Por tanto:
\[
(y_1,y_2,y_3)=(\alpha+\beta,\alpha-\beta,\beta).
\]

De aquí:
\[
y_3=\beta.
\]

Además:
\[
y_1-y_2=(\alpha+\beta)-(\alpha-\beta)=2\beta.
\]

Como \(\beta=y_3\), queda:
\[
y_1-y_2=2y_3.
\]

Luego:
\[
y_1-y_2-2y_3=0.
\]

Por tanto:
\[
\boxed{
\operatorname{Im} f=
\{(y_1,y_2,y_3)\in\mathbb{R}^3:y_1-y_2-2y_3=0\}.
}
\]

\subsubsection*{b) Clasificación de la aplicación lineal}

Como
\[
\ker f=L<\{(1,-1,0)\}>,
\]
tenemos:
\[
\ker f\neq \{(0,0,0)\}.
\]

Por tanto, \(f\) no es inyectiva.

Además:
\[
\dim(\operatorname{Im} f)=2.
\]

Como el espacio de llegada es \(\mathbb{R}^3\), se tiene:
\[
\dim(\operatorname{Im} f)=2<3=\dim(\mathbb{R}^3).
\]

Por tanto, \(f\) no es sobreyectiva.

En consecuencia:
\[
\boxed{
f \text{ no es inyectiva ni sobreyectiva.}
}
\]

Como \(f\) es un endomorfismo de \(\mathbb{R}^3\), y no es biyectiva, tampoco es un automorfismo.

\[
\boxed{
f \text{ no es un automorfismo.}
}
\]

\subsubsection*{c) Cálculo de \(f(V)\)}

Tenemos el subespacio
\[
V=\{(x_1,x_2,x_3)\in\mathbb{R}^3:x_1+x_2+x_3=0\}.
\]

Queremos calcular:
\[
f(V).
\]

Tomamos un vector cualquiera de \(V\):
\[
(x_1,x_2,x_3)\in V.
\]

Entonces cumple:
\[
x_1+x_2+x_3=0.
\]

Aplicamos \(f\):
\[
f(x_1,x_2,x_3)
=
(x_1+x_2+x_3,\;x_1+x_2-x_3,\;x_3).
\]

Como
\[
x_1+x_2+x_3=0,
\]
la primera coordenada de la imagen es:
\[
0.
\]

Además, de
\[
x_1+x_2+x_3=0
\]
obtenemos:
\[
x_1+x_2=-x_3.
\]

Por tanto:
\[
x_1+x_2-x_3=-x_3-x_3=-2x_3.
\]

Luego:
\[
f(x_1,x_2,x_3)=(0,-2x_3,x_3).
\]

Si llamamos
\[
x_3=t,
\]
entonces:
\[
f(x_1,x_2,x_3)=(0,-2t,t).
\]

Por tanto:
\[
f(V)=\{(0,-2t,t):t\in\mathbb{R}\}.
\]

Equivalentemente:
\[
f(V)=L<\{(0,-2,1)\}>.
\]

Así:
\[
\boxed{
f(V)=L<\{(0,-2,1)\}>.
}
\]

En forma implícita:
\[
\boxed{
f(V)=\{(y_1,y_2,y_3)\in\mathbb{R}^3:y_1=0,\ y_2+2y_3=0\}.
}
\]

Por tanto:
\[
\boxed{\dim f(V)=1.}
\]

\section*{Solución Ejercicio 2}

Tenemos la aplicación lineal
\[
f:\mathbb{R}^4\longrightarrow \mathbb{R}^3
\]
definida por
\[
f(x_1,x_2,x_3,x_4)=(x_1+x_2,\;0,\;x_1-x_2).
\]

La matriz asociada a \(f\) en las bases canónicas es
\[
M(f)=
\begin{pmatrix}
1&1&0&0\\
0&0&0&0\\
1&-1&0&0
\end{pmatrix}.
\]

\subsubsection*{Núcleo de \(f\)}

Para calcular el núcleo, resolvemos:
\[
f(x_1,x_2,x_3,x_4)=(0,0,0).
\]

Es decir:
\[
(x_1+x_2,\;0,\;x_1-x_2)=(0,0,0).
\]

Por tanto:
\[
\begin{cases}
x_1+x_2=0,\\
x_1-x_2=0.
\end{cases}
\]

Sumando ambas ecuaciones:
\[
2x_1=0,
\]
luego
\[
x_1=0.
\]

Sustituyendo en
\[
x_1+x_2=0,
\]
obtenemos:
\[
x_2=0.
\]

Las variables \(x_3\) y \(x_4\) son libres. Por tanto:
\[
x_3=\lambda,
\qquad
x_4=\mu,
\]
con \(\lambda,\mu\in\mathbb{R}\).

Así:
\[
(x_1,x_2,x_3,x_4)=(0,0,\lambda,\mu).
\]

Es decir:
\[
(x_1,x_2,x_3,x_4)
=
\lambda(0,0,1,0)+\mu(0,0,0,1).
\]

Por tanto,
\[
\ker f=L<\{(0,0,1,0),(0,0,0,1)\}>.
\]

Una base de \(\ker f\) es
\[
\boxed{\{(0,0,1,0),(0,0,0,1)\}.}
\]

Luego:
\[
\boxed{\dim(\ker f)=2.}
\]

\subsubsection*{Imagen de \(f\)}

A partir de la matriz:
\[
M(f)=
\begin{pmatrix}
1&1&0&0\\
0&0&0&0\\
1&-1&0&0
\end{pmatrix},
\]
la imagen está generada por sus columnas:
\[
\operatorname{Im} f
=
L<\{(1,0,1),(1,0,-1),(0,0,0),(0,0,0)\}>.
\]

Por tanto:
\[
\operatorname{Im} f
=
L<\{(1,0,1),(1,0,-1)\}>.
\]

Comprobamos que estos dos vectores son linealmente independientes.

Sea
\[
\alpha(1,0,1)+\beta(1,0,-1)=(0,0,0).
\]

Entonces:
\[
(\alpha+\beta,0,\alpha-\beta)=(0,0,0).
\]

Por tanto:
\[
\begin{cases}
\alpha+\beta=0,\\
\alpha-\beta=0.
\end{cases}
\]

Sumando ambas ecuaciones:
\[
2\alpha=0,
\]
luego
\[
\alpha=0.
\]

Sustituyendo:
\[
\beta=0.
\]

Luego son linealmente independientes.

Así, una base de la imagen es
\[
\boxed{\{(1,0,1),(1,0,-1)\}.}
\]

Por tanto:
\[
\boxed{\dim(\operatorname{Im} f)=2.}
\]

También podemos dar la imagen en forma implícita.

Como la segunda coordenada de cualquier vector imagen es siempre cero, se tiene:
\[
\operatorname{Im} f\subseteq \{(y_1,y_2,y_3)\in\mathbb{R}^3:y_2=0\}.
\]

Además, dado cualquier vector de la forma
\[
(y_1,0,y_3),
\]
queremos encontrar \(x_1,x_2\) tales que:
\[
x_1+x_2=y_1,
\qquad
x_1-x_2=y_3.
\]

Sumando:
\[
2x_1=y_1+y_3,
\]
luego:
\[
x_1=\frac{y_1+y_3}{2}.
\]

Restando:
\[
2x_2=y_1-y_3,
\]
luego:
\[
x_2=\frac{y_1-y_3}{2}.
\]

Por tanto, todo vector \((y_1,0,y_3)\) pertenece a la imagen.

Luego:
\[
\boxed{
\operatorname{Im} f=\{(y_1,y_2,y_3)\in\mathbb{R}^3:y_2=0\}.
}
\]

\subsubsection*{Comprobación con el teorema de la dimensión}

Como
\[
\dim(\mathbb{R}^4)=4,
\]
y
\[
\dim(\ker f)=2,
\]
por el teorema de la dimensión:
\[
\dim(\mathbb{R}^4)=\dim(\ker f)+\dim(\operatorname{Im} f).
\]

Entonces:
\[
4=2+\dim(\operatorname{Im} f).
\]

Por tanto:
\[
\dim(\operatorname{Im} f)=2,
\]
coincidiendo con lo obtenido anteriormente.

\section*{Solución Ejercicio 3}

Como trabajamos en \(P_2(x)\), aplicamos el isomorfismo natural
\[
\varphi:P_2(x)\longrightarrow \mathbb{R}^3,
\qquad
\varphi(a+bx+cx^2)=(a,b,c).
\]

Por tanto, identificamos cada polinomio
\[
a+bx+cx^2
\]
con el vector
\[
(a,b,c)\in\mathbb{R}^3.
\]

Bajo este isomorfismo, la base canónica de \(P_2(x)\),
\[
B_C=\{1,x,x^2\},
\]
se corresponde con la base canónica de \(\mathbb{R}^3\):
\[
\{(1,0,0),(0,1,0),(0,0,1)\}.
\]

Además,
\[
\ker f=\{(a,b,c)\in\mathbb{R}^3:a=b=c\}.
\]
Luego
\[
\ker f=L<\{(1,1,1)\}>.
\]

Por otra parte, los vectores que tienen término independiente nulo son los vectores de la forma
\[
(0,b,c).
\]
Como estos se transforman en sí mismos, se tiene:
\[
f(0,b,c)=(0,b,c).
\]

En particular,
\[
f(0,1,0)=(0,1,0),
\qquad
f(0,0,1)=(0,0,1).
\]

Es decir,
\[
f(e_2)=e_2,
\qquad
f(e_3)=e_3.
\]

\subsubsection*{a) Matriz de \(f\) en la base canónica}

Sabemos que
\[
(1,1,1)\in \ker f.
\]
Por tanto,
\[
f(1,1,1)=(0,0,0).
\]

Como
\[
(1,1,1)=e_1+e_2+e_3,
\]
por linealidad:
\[
f(e_1)+f(e_2)+f(e_3)=(0,0,0).
\]

Sustituyendo \(f(e_2)=e_2\) y \(f(e_3)=e_3\), queda:
\[
f(e_1)+e_2+e_3=(0,0,0).
\]

Luego
\[
f(e_1)=-e_2-e_3.
\]

Por tanto,
\[
f(e_1)=(0,-1,-1),
\qquad
f(e_2)=(0,1,0),
\qquad
f(e_3)=(0,0,1).
\]

La matriz de \(f\) en la base canónica se obtiene colocando estos vectores como columnas:
\[
M_{B_C}(f)=
\begin{pmatrix}
0&0&0\\
-1&1&0\\
-1&0&1
\end{pmatrix}.
\]

Por tanto,
\[
\boxed{
M_{B_C}(f)=
\begin{pmatrix}
0&0&0\\
-1&1&0\\
-1&0&1
\end{pmatrix}
}
\]

\subsubsection*{b) Expresión implícita de \(\operatorname{Im} f\)}

A partir de la matriz obtenida, si
\[
v=(a,b,c)\in\mathbb{R}^3,
\]
entonces
\[
f(a,b,c)=
\begin{pmatrix}
0&0&0\\
-1&1&0\\
-1&0&1
\end{pmatrix}
\begin{pmatrix}
a\\
b\\
c
\end{pmatrix}
=
\begin{pmatrix}
0\\
-a+b\\
-a+c
\end{pmatrix}.
\]

Luego todo vector de la imagen tiene primera coordenada nula. Por tanto,
\[
\operatorname{Im} f\subseteq \{(x,y,z)\in\mathbb{R}^3:x=0\}.
\]

Además,
\[
f(0,1,0)=(0,1,0),
\qquad
f(0,0,1)=(0,0,1).
\]
Luego los vectores \((0,1,0)\) y \((0,0,1)\) pertenecen a \(\operatorname{Im} f\), y generan todo el plano \(x=0\).

Así,
\[
\operatorname{Im} f=L<\{(0,1,0),(0,0,1)\}>.
\]

En forma implícita:
\[
\boxed{
\operatorname{Im} f=\{(x,y,z)\in\mathbb{R}^3:x=0\}
}
\]

Por tanto,
\[
\boxed{\dim(\operatorname{Im} f)=2.}
\]

\paragraph{Primera forma de razonar la dimensión.}

Como
\[
\operatorname{Im} f=\{(x,y,z)\in\mathbb{R}^3:x=0\},
\]
se trata de un subespacio de \(\mathbb{R}^3\) definido por una única ecuación lineal independiente. Por tanto,
\[
\dim(\operatorname{Im} f)=3-1=2.
\]

\paragraph{Segunda forma de razonar la dimensión.}

Por el teorema de la dimensión:
\[
\dim P_2(x)=\dim(\ker f)+\dim(\operatorname{Im} f).
\]

Como
\[
\dim P_2(x)=3
\]
y
\[
\ker f=L<\{(1,1,1)\}>,
\]
se tiene
\[
\dim(\ker f)=1.
\]

Por tanto,
\[
3=1+\dim(\operatorname{Im} f),
\]
de donde
\[
\dim(\operatorname{Im} f)=2.
\]

\subsubsection*{c) Cálculo de \(f(S)\)}

Tenemos
\[
S=\{(\lambda-2\mu+\delta,\lambda-2\mu-\delta,\lambda-2\mu): \lambda,\mu,\delta\in\mathbb{R}\}.
\]

Tomamos un vector cualquiera de \(S\):
\[
v=(\lambda-2\mu+\delta,\lambda-2\mu-\delta,\lambda-2\mu).
\]

Es decir,
\[
a=\lambda-2\mu+\delta,
\qquad
b=\lambda-2\mu-\delta,
\qquad
c=\lambda-2\mu.
\]

Como
\[
f(a,b,c)=(0,-a+b,-a+c),
\]
sustituimos:
\[
f(v)=
\left(
0,
-(\lambda-2\mu+\delta)+(\lambda-2\mu-\delta),
-(\lambda-2\mu+\delta)+(\lambda-2\mu)
\right).
\]

Simplificando:
\[
f(v)=(0,-2\delta,-\delta).
\]

Por tanto,
\[
f(S)=\{(0,-2\delta,-\delta):\delta\in\mathbb{R}\}.
\]

Equivalentemente,
\[
f(S)=L<\{(0,2,1)\}>.
\]

Luego
\[
\dim f(S)=1.
\]

Para obtener la forma implícita, si
\[
(x,y,z)\in f(S),
\]
entonces
\[
x=0
\]
y, como
\[
y=-2\delta,
\qquad
z=-\delta,
\]
se verifica
\[
y=2z.
\]

Por tanto,
\[
y-2z=0.
\]

Así,
\[
\boxed{
f(S)=\{(x,y,z)\in\mathbb{R}^3:x=0,\ y-2z=0\}
}
\]

\paragraph{¿Tiene sentido la dimensión obtenida?}

Sí tiene sentido.

Primero estudiamos \(S\). Escribimos
\[
(\lambda-2\mu+\delta,\lambda-2\mu-\delta,\lambda-2\mu)
=
(\lambda-2\mu)(1,1,1)+\delta(1,-1,0).
\]

Por tanto,
\[
S=L<\{(1,1,1),(1,-1,0)\}>.
\]

Los vectores \((1,1,1)\) y \((1,-1,0)\) son linealmente independientes, luego
\[
\dim S=2.
\]

Pero
\[
(1,1,1)\in \ker f.
\]

Esto significa que una de las direcciones de \(S\) desaparece al aplicar \(f\). En concreto,
\[
f(1,1,1)=(0,0,0).
\]

Por tanto, es coherente que \(S\) tenga dimensión \(2\), pero \(f(S)\) tenga dimensión \(1\).

De hecho,
\[
S\cap \ker f=L<\{(1,1,1)\}>,
\]
así que
\[
\dim f(S)=\dim S-\dim(S\cap \ker f)=2-1=1.
\]

\subsubsection*{d) Matriz de \(f\) en la base \(B^*\)}

La base dada es
\[
B^*=\{1+x+x^2,\;1+x,\;1\}.
\]

Bajo el isomorfismo
\[
P_2(x)\simeq \mathbb{R}^3,
\]
esta base se corresponde con
\[
B^*=\{(1,1,1),(1,1,0),(1,0,0)\}.
\]

Llamamos
\[
u_1=(1,1,1),
\qquad
u_2=(1,1,0),
\qquad
u_3=(1,0,0).
\]

\paragraph{Cálculo por coordenadas.}

Calculamos las imágenes de los vectores de la base \(B^*\).

Como
\[
u_1=(1,1,1)\in \ker f,
\]
se tiene
\[
f(u_1)=(0,0,0).
\]

Luego
\[
[f(u_1)]_{B^*}=
\begin{pmatrix}
0\\
0\\
0
\end{pmatrix}.
\]

Ahora calculamos \(f(u_2)\):
\[
f(u_2)=f(1,1,0)=(0,-1+1,-1+0)=(0,0,-1).
\]

Buscamos \(\alpha,\beta,\gamma\in\mathbb{R}\) tales que
\[
(0,0,-1)=\alpha(1,1,1)+\beta(1,1,0)+\gamma(1,0,0).
\]

Entonces
\[
(0,0,-1)=(\alpha+\beta+\gamma,\alpha+\beta,\alpha).
\]

Igualando coordenadas:
\[
\alpha=-1,
\]
\[
\alpha+\beta=0,
\]
\[
\alpha+\beta+\gamma=0.
\]

De aquí:
\[
\alpha=-1,\qquad \beta=1,\qquad \gamma=0.
\]

Por tanto,
\[
[f(u_2)]_{B^*}=
\begin{pmatrix}
-1\\
1\\
0
\end{pmatrix}.
\]

Ahora calculamos \(f(u_3)\):
\[
f(u_3)=f(1,0,0)=(0,-1,-1).
\]

Buscamos \(\alpha,\beta,\gamma\in\mathbb{R}\) tales que
\[
(0,-1,-1)=\alpha(1,1,1)+\beta(1,1,0)+\gamma(1,0,0).
\]

Entonces
\[
(0,-1,-1)=(\alpha+\beta+\gamma,\alpha+\beta,\alpha).
\]

Igualando coordenadas:
\[
\alpha=-1,
\]
\[
\alpha+\beta=-1,
\]
\[
\alpha+\beta+\gamma=0.
\]

De aquí:
\[
\alpha=-1,\qquad \beta=0,\qquad \gamma=1.
\]

Por tanto,
\[
[f(u_3)]_{B^*}=
\begin{pmatrix}
-1\\
0\\
1
\end{pmatrix}.
\]

La matriz \(M_{B^*}(f)\) se obtiene colocando estas coordenadas como columnas:
\[
M_{B^*}(f)=
\begin{pmatrix}
0&-1&-1\\
0&1&0\\
0&0&1
\end{pmatrix}.
\]

Por tanto,
\[
\boxed{
M_{B^*}(f)=
\begin{pmatrix}
0&-1&-1\\
0&1&0\\
0&0&1
\end{pmatrix}
}
\]

\paragraph{Cálculo matricial.}

La matriz de cambio de base de \(B^*\) a la base canónica \(B_C\) se construye colocando como columnas los vectores de \(B^*\) escritos en coordenadas canónicas:
\[
P_{B_C\leftarrow B^*}
=
\begin{pmatrix}
1&1&1\\
1&1&0\\
1&0&0
\end{pmatrix}.
\]

La matriz de \(f\) en la base canónica es
\[
M_{B_C}(f)=
\begin{pmatrix}
0&0&0\\
-1&1&0\\
-1&0&1
\end{pmatrix}.
\]

Por la fórmula de cambio de base:
\[
M_{B^*}(f)
=
P_{B_C\leftarrow B^*}^{-1}
M_{B_C}(f)
P_{B_C\leftarrow B^*}.
\]

Calculamos primero:
\[
M_{B_C}(f)P_{B_C\leftarrow B^*}
=
\begin{pmatrix}
0&0&0\\
-1&1&0\\
-1&0&1
\end{pmatrix}
\begin{pmatrix}
1&1&1\\
1&1&0\\
1&0&0
\end{pmatrix}
=
\begin{pmatrix}
0&0&0\\
0&0&-1\\
0&-1&-1
\end{pmatrix}.
\]

Además,
\[
P_{B_C\leftarrow B^*}^{-1}
=
\begin{pmatrix}
0&0&1\\
0&1&-1\\
1&-1&0
\end{pmatrix}.
\]

Por tanto,
\[
M_{B^*}(f)
=
\begin{pmatrix}
0&0&1\\
0&1&-1\\
1&-1&0
\end{pmatrix}
\begin{pmatrix}
0&0&0\\
0&0&-1\\
0&-1&-1
\end{pmatrix}
=
\begin{pmatrix}
0&-1&-1\\
0&1&0\\
0&0&1
\end{pmatrix}.
\]

Así, obtenemos de nuevo:
\[
\boxed{
M_{B^*}(f)=
\begin{pmatrix}
0&-1&-1\\
0&1&0\\
0&0&1
\end{pmatrix}.
}
\]

\subsubsection*{Conclusión}

La matriz de \(f\) en la base canónica es
\[
M_{B_C}(f)=
\begin{pmatrix}
0&0&0\\
-1&1&0\\
-1&0&1
\end{pmatrix}.
\]

La imagen de \(f\) es
\[
\operatorname{Im} f=\{(x,y,z)\in\mathbb{R}^3:x=0\},
\]
y
\[
\dim(\operatorname{Im} f)=2.
\]

Además,
\[
f(S)=\{(x,y,z)\in\mathbb{R}^3:x=0,\ y-2z=0\},
\]
con
\[
\dim f(S)=1.
\]

Finalmente, en la base
\[
B^*=\{1+x+x^2,\;1+x,\;1\},
\]
la matriz del endomorfismo es
\[
M_{B^*}(f)=
\begin{pmatrix}
0&-1&-1\\
0&1&0\\
0&0&1
\end{pmatrix}.
\]

\section*{Solución Ejercicio 4}

Tenemos las aplicaciones lineales
\[
f:\mathbb{R}^2\longrightarrow \mathbb{R}^3,
\qquad
g:\mathbb{R}^3\longrightarrow \mathbb{R}^2.
\]

Como nos dan las imágenes de los vectores de la base canónica, podemos construir directamente sus matrices asociadas.

\subsubsection*{Cálculo matricial}

La matriz de \(f\) se obtiene colocando como columnas las imágenes de los vectores de la base canónica de \(\mathbb{R}^2\):
\[
M(f)=
\begin{pmatrix}
2&3\\
1&2\\
2&1
\end{pmatrix}.
\]

La matriz de \(g\) se obtiene colocando como columnas las imágenes de los vectores de la base canónica de \(\mathbb{R}^3\):
\[
M(g)=
\begin{pmatrix}
3&1&1\\
-1&0&-2
\end{pmatrix}.
\]

Como
\[
h=g\circ f,
\]
entonces la matriz de \(h\) es
\[
M(h)=M(g)M(f).
\]

Calculamos:
\[
M(h)=
\begin{pmatrix}
3&1&1\\
-1&0&-2
\end{pmatrix}
\begin{pmatrix}
2&3\\
1&2\\
2&1
\end{pmatrix}.
\]

Por tanto,
\[
M(h)=
\begin{pmatrix}
3\cdot 2+1\cdot 1+1\cdot 2 & 3\cdot 3+1\cdot 2+1\cdot 1\\
-1\cdot 2+0\cdot 1-2\cdot 2 & -1\cdot 3+0\cdot 2-2\cdot 1
\end{pmatrix}.
\]

Luego
\[
M(h)=
\begin{pmatrix}
9&12\\
-6&-5
\end{pmatrix}.
\]

Ahora calculamos
\[
h(4,-2)=M(h)
\begin{pmatrix}
4\\
-2
\end{pmatrix}.
\]

Así,
\[
h(4,-2)=
\begin{pmatrix}
9&12\\
-6&-5
\end{pmatrix}
\begin{pmatrix}
4\\
-2
\end{pmatrix}
=
\begin{pmatrix}
9\cdot 4+12\cdot(-2)\\
-6\cdot 4+(-5)\cdot(-2)
\end{pmatrix}.
\]

Por tanto,
\[
h(4,-2)=
\begin{pmatrix}
12\\
-14
\end{pmatrix}.
\]

Luego
\[
\boxed{h(4,-2)=(12,-14).}
\]

\subsubsection*{Cálculo por coordenadas}

Como
\[
(4,-2)=4(1,0)-2(0,1),
\]
por linealidad de \(f\), tenemos:
\[
f(4,-2)=4f(1,0)-2f(0,1).
\]

Sustituimos:
\[
f(4,-2)=4(2,1,2)-2(3,2,1).
\]

Calculamos:
\[
f(4,-2)=(8,4,8)-(6,4,2).
\]

Luego
\[
f(4,-2)=(2,0,6).
\]

Ahora aplicamos \(g\):
\[
h(4,-2)=g(f(4,-2))=g(2,0,6).
\]

Como
\[
(2,0,6)=2(1,0,0)+0(0,1,0)+6(0,0,1),
\]
por linealidad de \(g\):
\[
g(2,0,6)=2g(1,0,0)+0g(0,1,0)+6g(0,0,1).
\]

Sustituimos:
\[
g(2,0,6)=2(3,-1)+0(1,0)+6(1,-2).
\]

Calculamos:
\[
g(2,0,6)=(6,-2)+(0,0)+(6,-12).
\]

Por tanto,
\[
g(2,0,6)=(12,-14).
\]

Así,
\[
\boxed{h(4,-2)=(12,-14).}
\]

\section*{Solución Ejercicio 5}

Tenemos las aplicaciones lineales
\[
f:\mathbb{R}^3\longrightarrow \mathbb{R}^2,
\qquad
f(x,y,z)=(2x-y,2y+z),
\]
y
\[
g:\mathbb{R}^2\longrightarrow \mathbb{R}^3,
\qquad
g(x,y)=(4x+2y,y,x+y).
\]

\subsubsection*{a) Matrices en las bases canónicas}

Empezamos calculando la matriz de \(f\) respecto de las bases canónicas.

Como
\[
f(x,y,z)=(2x-y,2y+z),
\]
tenemos:
\[
f(1,0,0)=(2,0),
\]
\[
f(0,1,0)=(-1,2),
\]
\[
f(0,0,1)=(0,1).
\]

Por tanto,
\[
M_{CC'}(f)=
\begin{pmatrix}
2&-1&0\\
0&2&1
\end{pmatrix}.
\]

Ahora calculamos la matriz de \(g\) respecto de las bases canónicas.

Como
\[
g(x,y)=(4x+2y,y,x+y),
\]
tenemos:
\[
g(1,0)=(4,0,1),
\]
\[
g(0,1)=(2,1,1).
\]

Por tanto,
\[
M_{C'C}(g)=
\begin{pmatrix}
4&2\\
0&1\\
1&1
\end{pmatrix}.
\]

Así,
\[
\boxed{
M_{CC'}(f)=
\begin{pmatrix}
2&-1&0\\
0&2&1
\end{pmatrix}
}
\]
y
\[
\boxed{
M_{C'C}(g)=
\begin{pmatrix}
4&2\\
0&1\\
1&1
\end{pmatrix}.
}
\]

\subsubsection*{b) Base de \(\ker f\), base de \(\operatorname{Im} g\) y clasificación}

Para calcular el núcleo de \(f\), resolvemos:
\[
f(x,y,z)=(0,0).
\]

Es decir,
\[
(2x-y,2y+z)=(0,0).
\]

Por tanto:
\[
2x-y=0,
\]
\[
2y+z=0.
\]

De la primera ecuación:
\[
y=2x.
\]

Sustituyendo en la segunda:
\[
2(2x)+z=0,
\]
\[
4x+z=0,
\]
\[
z=-4x.
\]

Luego:
\[
(x,y,z)=(x,2x,-4x)=x(1,2,-4).
\]

Por tanto,
\[
\ker f=L<\{(1,2,-4)\}>.
\]

Una base de \(\ker f\) es
\[
\boxed{\{(1,2,-4)\}.}
\]

En consecuencia,
\[
\dim(\ker f)=1.
\]

Como
\[
\dim(\mathbb{R}^3)=3,
\]
por el teorema de la dimensión:
\[
\dim(\mathbb{R}^3)=\dim(\ker f)+\dim(\operatorname{Im} f).
\]

Entonces:
\[
3=1+\dim(\operatorname{Im} f),
\]
de donde
\[
\dim(\operatorname{Im} f)=2.
\]

Como \(f:\mathbb{R}^3\longrightarrow \mathbb{R}^2\), y
\[
\dim(\operatorname{Im} f)=2=\dim(\mathbb{R}^2),
\]
deducimos que \(f\) es sobreyectiva.

Pero como
\[
\ker f\neq \{(0,0,0)\},
\]
la aplicación \(f\) no es inyectiva.

Por tanto:
\[
\boxed{f \text{ es sobreyectiva, pero no inyectiva.}}
\]

Ahora calculamos la imagen de \(g\).

Sabemos que
\[
g(1,0)=(4,0,1),
\]
\[
g(0,1)=(2,1,1).
\]

Luego
\[
\operatorname{Im} g=L<\{(4,0,1),(2,1,1)\}>.
\]

Comprobamos si estos dos vectores son linealmente independientes.

Supongamos que
\[
\alpha(4,0,1)+\beta(2,1,1)=(0,0,0).
\]

Entonces:
\[
(4\alpha+2\beta,\beta,\alpha+\beta)=(0,0,0).
\]

De la segunda coordenada:
\[
\beta=0.
\]

Sustituyendo en la tercera:
\[
\alpha+\beta=0,
\]
\[
\alpha=0.
\]

Por tanto, son linealmente independientes.

Una base de \(\operatorname{Im} g\) es
\[
\boxed{\{(4,0,1),(2,1,1)\}.}
\]

Además,
\[
\dim(\operatorname{Im} g)=2.
\]

Como
\[
g:\mathbb{R}^2\longrightarrow \mathbb{R}^3
\]
y
\[
\dim(\operatorname{Im} g)=2=\dim(\mathbb{R}^2),
\]
deducimos que
\[
\dim(\ker g)=0.
\]

Por tanto,
\[
\ker g=\{(0,0)\}.
\]

Luego \(g\) es inyectiva.

Sin embargo,
\[
\dim(\operatorname{Im} g)=2<3=\dim(\mathbb{R}^3),
\]
así que \(g\) no es sobreyectiva.

Por tanto:
\[
\boxed{g \text{ es inyectiva, pero no sobreyectiva.}}
\]

\subsubsection*{c) Cálculo de \([f(2,0,-1)]_{B'}\) utilizando \(M_{BB'}(f)\)}

Queremos calcular
\[
[f(2,0,-1)]_{B'}
\]
utilizando la matriz de \(f\) respecto de la base \(B\) en el dominio y la base \(B'\) en el codominio.

Recordamos que
\[
B=\{(1,0,0),(1,1,0),(1,1,1)\}.
\]

Llamamos:
\[
u_1=(1,0,0),
\qquad
u_2=(1,1,0),
\qquad
u_3=(1,1,1).
\]

Y en \(\mathbb{R}^2\):
\[
B'=\{(1,0),(1,1)\}.
\]

Llamamos:
\[
v_1=(1,0),
\qquad
v_2=(1,1).
\]

Primero calculamos las imágenes de los vectores de la base \(B\).

Para \(u_1=(1,0,0)\):
\[
f(u_1)=f(1,0,0)=(2,0).
\]

Buscamos sus coordenadas en la base \(B'\):
\[
(2,0)=\alpha(1,0)+\beta(1,1).
\]

Entonces:
\[
(2,0)=(\alpha+\beta,\beta).
\]

Igualando coordenadas:
\[
\beta=0,
\]
\[
\alpha+\beta=2.
\]

Luego:
\[
\alpha=2,
\qquad
\beta=0.
\]

Por tanto:
\[
[f(u_1)]_{B'}=
\begin{pmatrix}
2\\
0
\end{pmatrix}.
\]

Para \(u_2=(1,1,0)\):
\[
f(u_2)=f(1,1,0)=(1,2).
\]

Buscamos sus coordenadas en \(B'\):
\[
(1,2)=\alpha(1,0)+\beta(1,1).
\]

Entonces:
\[
(1,2)=(\alpha+\beta,\beta).
\]

Igualando coordenadas:
\[
\beta=2,
\]
\[
\alpha+\beta=1.
\]

Luego:
\[
\alpha=-1,
\qquad
\beta=2.
\]

Por tanto:
\[
[f(u_2)]_{B'}=
\begin{pmatrix}
-1\\
2
\end{pmatrix}.
\]

Para \(u_3=(1,1,1)\):
\[
f(u_3)=f(1,1,1)=(1,3).
\]

Buscamos sus coordenadas en \(B'\):
\[
(1,3)=\alpha(1,0)+\beta(1,1).
\]

Entonces:
\[
(1,3)=(\alpha+\beta,\beta).
\]

Igualando coordenadas:
\[
\beta=3,
\]
\[
\alpha+\beta=1.
\]

Luego:
\[
\alpha=-2,
\qquad
\beta=3.
\]

Por tanto:
\[
[f(u_3)]_{B'}=
\begin{pmatrix}
-2\\
3
\end{pmatrix}.
\]

Colocando estas coordenadas como columnas, obtenemos:
\[
M_{BB'}(f)=
\begin{pmatrix}
2&-1&-2\\
0&2&3
\end{pmatrix}.
\]

Ahora escribimos el vector \((2,0,-1)\) en coordenadas de la base \(B\).

Buscamos \(\alpha,\beta,\gamma\in\mathbb{R}\) tales que
\[
(2,0,-1)=\alpha(1,0,0)+\beta(1,1,0)+\gamma(1,1,1).
\]

Entonces:
\[
(2,0,-1)=(\alpha+\beta+\gamma,\beta+\gamma,\gamma).
\]

Igualando coordenadas:
\[
\gamma=-1,
\]
\[
\beta+\gamma=0,
\]
\[
\alpha+\beta+\gamma=2.
\]

De \(\beta+\gamma=0\), como \(\gamma=-1\), tenemos:
\[
\beta=1.
\]

Y de \(\alpha+\beta+\gamma=2\):
\[
\alpha+1-1=2,
\]
\[
\alpha=2.
\]

Por tanto:
\[
[(2,0,-1)]_B=
\begin{pmatrix}
2\\
1\\
-1
\end{pmatrix}.
\]

Ahora aplicamos la matriz \(M_{BB'}(f)\):
\[
[f(2,0,-1)]_{B'}
=
M_{BB'}(f)[(2,0,-1)]_B.
\]

Es decir:
\[
[f(2,0,-1)]_{B'}
=
\begin{pmatrix}
2&-1&-2\\
0&2&3
\end{pmatrix}
\begin{pmatrix}
2\\
1\\
-1
\end{pmatrix}.
\]

Calculamos:
\[
[f(2,0,-1)]_{B'}
=
\begin{pmatrix}
2\cdot 2+(-1)\cdot 1+(-2)\cdot(-1)\\
0\cdot 2+2\cdot 1+3\cdot(-1)
\end{pmatrix}.
\]

Luego:
\[
[f(2,0,-1)]_{B'}
=
\begin{pmatrix}
5\\
-1
\end{pmatrix}.
\]

Por tanto:
\[
\boxed{
[f(2,0,-1)]_{B'}=
\begin{pmatrix}
5\\
-1
\end{pmatrix}
}
\]

\subsubsection*{Comprobación en las bases canónicas}

Calculamos directamente:
\[
f(2,0,-1)=(2\cdot 2-0,2\cdot 0+(-1)).
\]

Por tanto:
\[
f(2,0,-1)=(4,-1).
\]

Ahora comprobamos que sus coordenadas en la base \(B'\) son las obtenidas anteriormente.

Buscamos \(\alpha,\beta\in\mathbb{R}\) tales que
\[
(4,-1)=\alpha(1,0)+\beta(1,1).
\]

Entonces:
\[
(4,-1)=(\alpha+\beta,\beta).
\]

Igualando coordenadas:
\[
\beta=-1,
\]
\[
\alpha+\beta=4.
\]

Como \(\beta=-1\), tenemos:
\[
\alpha-1=4,
\]
\[
\alpha=5.
\]

Por tanto:
\[
[(4,-1)]_{B'}=
\begin{pmatrix}
5\\
-1
\end{pmatrix}.
\]

Así:
\[
\boxed{
[f(2,0,-1)]_{B'}=
\begin{pmatrix}
5\\
-1
\end{pmatrix}
}
\]
como queríamos comprobar.

\subsubsection*{Cálculo matricial de \(M_{BB'}(f)\)}

También podemos obtener \(M_{BB'}(f)\) mediante matrices de cambio de base.

La matriz de cambio de base de \(B\) a la base canónica \(C\) es:
\[
P_{C\leftarrow B}
=
\begin{pmatrix}
1&1&1\\
0&1&1\\
0&0&1
\end{pmatrix}.
\]

La matriz de cambio de base de \(B'\) a la base canónica \(C'\) es:
\[
P_{C'\leftarrow B'}
=
\begin{pmatrix}
1&1\\
0&1
\end{pmatrix}.
\]

Por tanto:
\[
P_{B'\leftarrow C'}=
P_{C'\leftarrow B'}^{-1}
=
\begin{pmatrix}
1&-1\\
0&1
\end{pmatrix}.
\]

Sabemos que:
\[
M_{CC'}(f)=
\begin{pmatrix}
2&-1&0\\
0&2&1
\end{pmatrix}.
\]

Entonces:
\[
M_{BB'}(f)
=
P_{B'\leftarrow C'}M_{CC'}(f)P_{C\leftarrow B}.
\]

Sustituyendo:
\[
M_{BB'}(f)
=
\begin{pmatrix}
1&-1\\
0&1
\end{pmatrix}
\begin{pmatrix}
2&-1&0\\
0&2&1
\end{pmatrix}
\begin{pmatrix}
1&1&1\\
0&1&1\\
0&0&1
\end{pmatrix}.
\]

Primero:
\[
\begin{pmatrix}
2&-1&0\\
0&2&1
\end{pmatrix}
\begin{pmatrix}
1&1&1\\
0&1&1\\
0&0&1
\end{pmatrix}
=
\begin{pmatrix}
2&1&1\\
0&2&3
\end{pmatrix}.
\]

Luego:
\[
M_{BB'}(f)
=
\begin{pmatrix}
1&-1\\
0&1
\end{pmatrix}
\begin{pmatrix}
2&1&1\\
0&2&3
\end{pmatrix}
=
\begin{pmatrix}
2&-1&-2\\
0&2&3
\end{pmatrix}.
\]

Así:
\[
\boxed{
M_{BB'}(f)=
\begin{pmatrix}
2&-1&-2\\
0&2&3
\end{pmatrix}.
}
\]

\section*{Solución Ejercicio 6}

Tenemos una aplicación lineal
\[
f:\mathbb{R}^3\longrightarrow \mathbb{R}^2
\]
tal que
\[
\ker f=\{(x,y,z)\in\mathbb{R}^3:x+ay+z=0,\ ax+y+z=0\}.
\]

El núcleo viene dado por el sistema homogéneo:
\[
\begin{cases}
x+ay+z=0,\\
ax+y+z=0.
\end{cases}
\]

\subsubsection*{a) Dimensiones y bases de \(\ker f\) y de \(\operatorname{Im} f\)}

Estudiamos el sistema que define el núcleo:
\[
\begin{cases}
x+ay+z=0,\\
ax+y+z=0.
\end{cases}
\]

Restamos ambas ecuaciones:
\[
(x+ay+z)-(ax+y+z)=0.
\]

Entonces:
\[
(1-a)x+(a-1)y=0.
\]

Sacamos factor común:
\[
(1-a)(x-y)=0.
\]

Distinguimos dos casos.

\paragraph{Caso 1: \(a\neq 1\).}

Si
\[
a\neq 1,
\]
entonces
\[
1-a\neq 0,
\]
y por tanto:
\[
x-y=0.
\]

Luego:
\[
x=y.
\]

Sustituimos en la primera ecuación:
\[
x+ax+z=0.
\]

Por tanto:
\[
(a+1)x+z=0,
\]
de donde:
\[
z=-(a+1)x.
\]

Así, los vectores del núcleo son:
\[
(x,y,z)=(x,x,-(a+1)x).
\]

Es decir:
\[
(x,y,z)=x(1,1,-a-1).
\]

Por tanto:
\[
\ker f=L<\{(1,1,-a-1)\}>.
\]

Una base de \(\ker f\) es:
\[
\boxed{\{(1,1,-a-1)\}.}
\]

Luego:
\[
\boxed{\dim(\ker f)=1.}
\]

Por el teorema de la dimensión:
\[
\dim(\mathbb{R}^3)=\dim(\ker f)+\dim(\operatorname{Im} f).
\]

Como
\[
\dim(\mathbb{R}^3)=3
\]
y
\[
\dim(\ker f)=1,
\]
tenemos:
\[
3=1+\dim(\operatorname{Im} f).
\]

Por tanto:
\[
\dim(\operatorname{Im} f)=2.
\]

Como
\[
f:\mathbb{R}^3\longrightarrow \mathbb{R}^2,
\]
se tiene:
\[
\operatorname{Im} f=\mathbb{R}^2.
\]

Luego una base de \(\operatorname{Im} f\) es:
\[
\boxed{\{(1,0),(0,1)\}.}
\]

\paragraph{Caso 2: \(a=1\).}

Si
\[
a=1,
\]
las dos ecuaciones que definen el núcleo coinciden:
\[
x+y+z=0.
\]

Por tanto:
\[
\ker f=\{(x,y,z)\in\mathbb{R}^3:x+y+z=0\}.
\]

Despejamos:
\[
x=-y-z.
\]

Tomamos
\[
y=\lambda,
\qquad
z=\mu.
\]

Entonces:
\[
(x,y,z)=(-\lambda-\mu,\lambda,\mu).
\]

Escribimos:
\[
(x,y,z)=\lambda(-1,1,0)+\mu(-1,0,1).
\]

Por tanto:
\[
\ker f=L<\{(-1,1,0),(-1,0,1)\}>.
\]

Una base de \(\ker f\) es:
\[
\boxed{\{(-1,1,0),(-1,0,1)\}.}
\]

Luego:
\[
\boxed{\dim(\ker f)=2.}
\]

Por el teorema de la dimensión:
\[
3=\dim(\ker f)+\dim(\operatorname{Im} f).
\]

Como
\[
\dim(\ker f)=2,
\]
tenemos:
\[
3=2+\dim(\operatorname{Im} f).
\]

Por tanto:
\[
\dim(\operatorname{Im} f)=1.
\]

Además, sabemos que:
\[
f(1,0,1)=(2,1).
\]

Como \((2,1)\neq (0,0)\), este vector genera la imagen, porque la imagen tiene dimensión \(1\).

Así:
\[
\operatorname{Im} f=L<\{(2,1)\}>.
\]

Una base de \(\operatorname{Im} f\) es:
\[
\boxed{\{(2,1)\}.}
\]

\subsubsection*{b) Matriz de \(f\) en las bases canónicas}

Queremos calcular la matriz de \(f\) en las bases canónicas de \(\mathbb{R}^3\) y \(\mathbb{R}^2\).

Distinguimos de nuevo dos casos.

\paragraph{Caso 1: \(a\neq 1\).}

Como
\[
\ker f=\{(x,y,z)\in\mathbb{R}^3:x+ay+z=0,\ ax+y+z=0\},
\]
las filas de la matriz de \(f\) deben ser combinaciones lineales de los vectores:
\[
(1,a,1),
\qquad
(a,1,1).
\]

Es decir, podemos escribir:
\[
f(x,y,z)=
\left(
\alpha(x+ay+z)+\beta(ax+y+z),
\gamma(x+ay+z)+\delta(ax+y+z)
\right).
\]

Usamos las condiciones dadas.

Primero:
\[
f(0,0,a-1)=(0,a-1).
\]

Calculamos:
\[
x+ay+z=a-1,
\qquad
ax+y+z=a-1.
\]

Por tanto:
\[
f(0,0,a-1)=
\left(
(\alpha+\beta)(a-1),
(\gamma+\delta)(a-1)
\right).
\]

Como \(a\neq 1\), tenemos \(a-1\neq 0\). Entonces:
\[
\alpha+\beta=0,
\]
\[
\gamma+\delta=1.
\]

Ahora usamos:
\[
f(1,0,1)=(2,1).
\]

Para \((1,0,1)\):
\[
x+ay+z=2,
\qquad
ax+y+z=a+1.
\]

Luego:
\[
f(1,0,1)=
\left(
2\alpha+(a+1)\beta,
2\gamma+(a+1)\delta
\right).
\]

Por tanto:
\[
2\alpha+(a+1)\beta=2,
\]
\[
2\gamma+(a+1)\delta=1.
\]

De
\[
\alpha+\beta=0
\]
tenemos:
\[
\beta=-\alpha.
\]

Sustituyendo:
\[
2\alpha-(a+1)\alpha=2.
\]

Entonces:
\[
(1-a)\alpha=2,
\]
y como \(a\neq 1\):
\[
\alpha=\frac{2}{1-a}.
\]

Luego:
\[
\beta=-\frac{2}{1-a}.
\]

La primera componente de \(f\) queda:
\[
\frac{2}{1-a}(x+ay+z)-\frac{2}{1-a}(ax+y+z).
\]

Simplificando:
\[
\frac{2}{1-a}\left((x+ay+z)-(ax+y+z)\right).
\]

\[
=\frac{2}{1-a}\left((1-a)x+(a-1)y\right).
\]

\[
=2x-2y.
\]

Por tanto, la primera fila de la matriz es:
\[
(2,-2,0).
\]

Ahora resolvemos la segunda componente.

De
\[
\gamma+\delta=1,
\]
tenemos:
\[
\gamma=1-\delta.
\]

Sustituyendo en
\[
2\gamma+(a+1)\delta=1,
\]
obtenemos:
\[
2(1-\delta)+(a+1)\delta=1.
\]

Entonces:
\[
2-2\delta+(a+1)\delta=1.
\]

\[
2+(a-1)\delta=1.
\]

Por tanto:
\[
(a-1)\delta=-1,
\]
y como \(a\neq 1\):
\[
\delta=-\frac{1}{a-1}.
\]

Entonces:
\[
\gamma=1+\frac{1}{a-1}
=
\frac{a}{a-1}.
\]

La segunda componente de \(f\) queda:
\[
\frac{a}{a-1}(x+ay+z)-\frac{1}{a-1}(ax+y+z).
\]

Simplificando:
\[
\frac{1}{a-1}
\left(
a(x+ay+z)-(ax+y+z)
\right).
\]

\[
=
\frac{1}{a-1}
\left(
ax+a^2y+az-ax-y-z
\right).
\]

\[
=
\frac{1}{a-1}
\left(
(a^2-1)y+(a-1)z
\right).
\]

\[
=(a+1)y+z.
\]

Por tanto, la segunda fila de la matriz es:
\[
(0,a+1,1).
\]

Luego, si \(a\neq 1\),
\[
M_{CC'}(f)=
\begin{pmatrix}
2&-2&0\\
0&a+1&1
\end{pmatrix}.
\]

Así:
\[
\boxed{
M_{CC'}(f)=
\begin{pmatrix}
2&-2&0\\
0&a+1&1
\end{pmatrix},
\qquad a\neq 1.
}
\]

\paragraph{Caso 2: \(a=1\).}

Si
\[
a=1,
\]
tenemos:
\[
\ker f=\{(x,y,z)\in\mathbb{R}^3:x+y+z=0\}.
\]

Por tanto, cada componente de \(f\) debe ser proporcional a:
\[
x+y+z.
\]

Escribimos:
\[
f(x,y,z)=(\alpha(x+y+z),\beta(x+y+z)).
\]

Usamos la condición:
\[
f(1,0,1)=(2,1).
\]

Como
\[
1+0+1=2,
\]
tenemos:
\[
f(1,0,1)=(2\alpha,2\beta).
\]

Igualando:
\[
(2\alpha,2\beta)=(2,1).
\]

Por tanto:
\[
\alpha=1,
\qquad
\beta=\frac{1}{2}.
\]

Así:
\[
f(x,y,z)=\left(x+y+z,\frac{1}{2}(x+y+z)\right).
\]

La matriz de \(f\) es:
\[
M_{CC'}(f)=
\begin{pmatrix}
1&1&1\\
\frac{1}{2}&\frac{1}{2}&\frac{1}{2}
\end{pmatrix}.
\]

Por tanto:
\[
\boxed{
M_{CC'}(f)=
\begin{pmatrix}
1&1&1\\
\frac{1}{2}&\frac{1}{2}&\frac{1}{2}
\end{pmatrix},
\qquad a=1.
}
\]

\subsubsection*{c) Para \(a=1\), hallar las pre imagenes de \((2,1)\) y \((2,3)\)}

Para \(a=1\), hemos obtenido:
\[
f(x,y,z)=\left(x+y+z,\frac{1}{2}(x+y+z)\right).
\]

Es decir, si llamamos
\[
s=x+y+z,
\]
entonces:
\[
f(x,y,z)=\left(s,\frac{s}{2}\right).
\]

\paragraph{Vectores que se transforman en \((2,1)\).}

Queremos resolver:
\[
f(x,y,z)=(2,1).
\]

Esto equivale a:
\[
\left(s,\frac{s}{2}\right)=(2,1).
\]

Por tanto:
\[
s=2
\]
y
\[
\frac{s}{2}=1.
\]

Ambas condiciones son compatibles, pues si \(s=2\), entonces:
\[
\frac{s}{2}=1.
\]

Luego necesitamos:
\[
x+y+z=2.
\]

Por tanto, los vectores que se transforman en \((2,1)\) son:
\[
\{(x,y,z)\in\mathbb{R}^3:x+y+z=2\}.
\]

Parametrizamos:
\[
y=\lambda,
\qquad
z=\mu.
\]

Entonces:
\[
x=2-\lambda-\mu.
\]

Así:
\[
(x,y,z)=(2-\lambda-\mu,\lambda,\mu),
\qquad
\lambda,\mu\in\mathbb{R}.
\]

Por tanto:
\[
\boxed{
f^{-1}(2,1)=\{(2-\lambda-\mu,\lambda,\mu):\lambda,\mu\in\mathbb{R}\}.
}
\]

\paragraph{Vectores que se transforman en \((2,3)\).}

Queremos resolver:
\[
f(x,y,z)=(2,3).
\]

Esto equivale a:
\[
\left(s,\frac{s}{2}\right)=(2,3).
\]

Por tanto:
\[
s=2
\]
y
\[
\frac{s}{2}=3.
\]

Pero si
\[
s=2,
\]
entonces:
\[
\frac{s}{2}=1,
\]
no \(3\).

Por tanto, el sistema es incompatible.

Luego no existe ningún vector de \(\mathbb{R}^3\) que se transforme en \((2,3)\).

Así:
\[
\boxed{
f^{-1}(2,3)=\varnothing.
}
\]

Esto también se puede razonar observando que, para \(a=1\),
\[
\operatorname{Im} f=L<\{(2,1)\}>.
\]

Como
\[
(2,1)\in \operatorname{Im} f,
\]
sí tiene pre imagen.

En cambio,
\[
(2,3)\notin L<\{(2,1)\}>,
\]
por lo que no tiene pre imagen.
\end{document}